\topic{A short first route, with real return visits}
\labelled{First visit.}{Allow the usual arrival/movement time, then about 5 minutes for handling and the shared crossing. Give pairs student p. 1 for 10--15 minutes, then p. 2 for 15--20 minutes if the rule is secure. Close with 5 minutes comparing two saved trips. One carefully checked trip and a conjecture can be a good stopping point. If time and attention remain, start p. 3; completing every page is not the aim.}
\labelled{Return visit: one straight developed line.}{Reestablish one crossing, then give p. 3 a full 15--25 minutes. Compare complete two-piece loops, move the start, and check how separately numbered copies make the line straight. This adds a representation for explaining the portal rule.}
\labelled{Return visit: the exceptional point.}{Use pp. 4--5 together. Allow 10--15 minutes for endpoint matching and another 15--20 for cutting, ordering and counting sectors. Keep an uncut p. 4 beside the pieces; precut p. 5 if cutting would use the whole session. This is a different question from path tracing.}
\labelled{Return visit: exact returns.}{Start the L model on p. 6, with its new square labels and outside edge pairs. Continue to p. 7 or p. 8 when children can check that rule. Save p. 9 and the universal explanation for a later visit if needed. Each deserves 15--25 minutes, not a final five-minute challenge. The p. 9 proof is a readiness-dependent Grades 4--5 extension.}
\labelled{If the rule is too much.}{Return to one crossing at a time, with a partner controlling the direction arrow and an adult steadying a tracing. If the adult must choose every point and direction, stop the independent investigation and use a familiar activity. This packet makes no K--1 entry claim.}
\topic{Student p. 1 / Problem 1: three marked trips}
\route{Concrete entry. All three marked rightward paths close; none hits a corner. Edge letters A, B, C, D refer to printed pairs, not the later L-square names.}
\begin{tabularx}{\linewidth}{@{}p{.40in}p{1.3in}X@{}}
\toprule
Start & Crossings in order & Full answer \\
\midrule
P & C & Stays in the middle band; one loop is the full left-to-right width. \\
Q & B, D & Runs through the upper and lower bands; two crossings before returning to Q. \\
R & D, B & Runs through lower and upper bands; two crossings before returning to R. \\
\bottomrule
\end{tabularx}\par
\labelled{Exact check.}{In the background coordinates, \(P=(-.65,0)\), \(Q=(-.3,1.60)\), \(R=(.20,-1.45)\). P has length \(2a\). Q and R each have length \(2b\), although their paired chords differ. Q's lower height is \(1.60-b\); R's upper height is \(b-1.45\). These are distinct loops in the same outer family.}
\hint{1. Have the checker identify the mate before moving. 2. Trace the edge and mark its crossing point, then slide that tracing to the paired edge without rotation. 3. When the route crosses its starting horizontal line, keep going until it reaches the actual starting dot.}
\topic{Student p. 2 / Problem 2: classify all interior starts}
\labelled{Complete answer.}{There are exactly two closed-trip lengths. Starts strictly between the horizontal lines through the vertical-edge endpoints give the shorter \(2a\) loop. Starts strictly above or below that central band give the longer \(2b\) loop. Interior starts on the two dividing lines hit a corner and stop. No other interior cases exist; the horizontal coordinate only changes where the loop begins.}
\labelled{Measurement.}{On this p. 2 octagon, printed width is 5.35 inches. At actual size, short loops measure 5.35 inches and long loops \(5.35\sqrt2\approx7.57\) inches. These are full traveled lengths, with jumps excluded and the final piece included. Never compare inches across different-sized pages.}
\hint{1. Compare a near-top start with a central one. 2. Trace both pieces of a complete outer trip onto one strip. 3. Move a start vertically and look for where its exit edge changes. The formula \(2(b-y)+2y\) explains why measurements that look nearly equal really are equal.}
\newpage
\topic{Student p. 3 / Problem 3: a balanced loop, drawn straight}
\route{First substantial extension or return visit. The printed task starts on an edge, away from a corner, so one trip consists of two complete chords.}
\labelled{Answer.}{The upper and lower chords must have equal length. In side-length-2 coordinates, their heights are \(y=\pm b/2=\pm(2+\sqrt2)/2\); each chord has length \(b\). Start at the midpoint of the upper-left D edge, heading right. Exit at the upper-right B midpoint, resume at the lower-left B midpoint, and cross the lower chord to the lower-right D midpoint. This is the matching point of the starting edge, with the same heading: one full trip. A lower-left B midpoint start gives the same loop in the other order.}
\labelled{Printed seam-start development.}{Use copy 1 at its original position and copy 2 translated by \((b,b)\), putting its SW B edge on copy 1's NE B edge. The straight line from \((-b/2,b/2)\) to \((3b/2,b/2)\) traverses the two full chords. Its final point, on copy 2's D edge, matches the original start. A third copy may mark that correspondence, but its interior is not needed for the completed trip.}
\labelled{Optional interior-start variant, shown below.}{After the printed task, move the start to \(S=(0,b/2)\). Keep copy 2 as above and translate copy 3 by \((2b,0)\), matching its NW D edge to copy 2's SE D edge. The line at \(y=b/2\) reaches the same local point \(S\) in copy 3 after distance \(2b\). This variant has three drawn pieces because one complete chord is split at the start.}
\begin{center}
\begin{tikzpicture}[scale=.54,>=Stealth,every node/.style={font=\small}]
\begin{scope}\octagon\node at (0,-.35){copy 1};\end{scope}
\begin{scope}[shift={(3.41421356,3.41421356)}]\octagon\node at (0,.55){copy 2};\end{scope}
\begin{scope}[shift={(6.82842712,0)}]\octagon\node at (0,-.35){copy 3};\end{scope}
\draw[pathblue,line width=1.4pt,->] (0,1.70710678)--(6.82842712,1.70710678);
\foreach \x in {0,1.70710678,5.12132034,6.82842712}{\fill[pathblue] (\x,1.70710678) circle(2.2pt);}
\node[below] at (0,1.52){start \(S\)};
\node[below] at (1.70710678,1.70710678){B};
\node[below] at (5.12132034,1.70710678){D};
\node[below] at (6.82842712,1.52){return \(S\)};
\end{tikzpicture}
\end{center}
\labelled{Why three copies in the optional picture?}{An interior start cuts the upper chord into an initial and final piece. Here their lengths are \(b/2+b+b/2=2b\). A different interior start changes the first and last lengths, keeping their sum \(b\). The printed seam-start route instead uses two whole chords. In either drawing the final point represents the starting surface point, even though it lies elsewhere on the desk.}
\hint{1. Compare the complete upper and lower chords, not the first piece from the dot. 2. Move the upper height until neither piece looks longer. 3. Try the midpoint of a sloping edge and check its mate. For developing, align two labeled edges first; then move the ruler.}
\labelled{Checks to ask for.}{Copies must not rotate or flip; the developed path is one line. Start and return must be matching surface points with the same arrow. The printed convention uses two whole chords; an independently correct interior-start drawing is a useful variant, not a failure. Do not demand a plane-filling pattern. If later copies overlap, keep their sheet identities distinct.}
\labelled{An explanation without algebra.}{Mark the upper and lower pieces on tracing strips. Move the upper line upward a little: both sloping ends move inward, while the partnered lower line expands by the same amount. At the balanced middle position the two strips match. The midpoint construction gives an exact check of an experimental guess.}
\newpage
\topic{Student p. 4 / Problem 4: which corners are one point?}
\route{Corner-identification investigation. Keep each polygon flat and trace its endpoint matches. Keep the uncut diagrams visible for the next problem.}
\labelled{Octagon answer.}{One group contains all eight corners: \(\{1,2,3,4,5,6,7,8\}\). The complete endpoint pairs, using printed labels, are
\[
\begin{array}{c|cc}
\text{edge pair}&\multicolumn{2}{c}{\text{endpoint matches}}\\ \hline
A&8\sim5&1\sim4\\
B&1\sim6&2\sim5\\
C&3\sim6&2\sim7\\
D&4\sim7&3\sim8.
\end{array}
\]
For a single checkable tour through the group, use \(1\to6\to3\to8\to5\to2\to7\to4\to1\). Every arrow is an allowed endpoint match, not a drawn path through the polygon's interior.}
\labelled{Square answer.}{One group contains all four corners: \(\{9,10,11,12\}\). E matches \(12\sim11\), \(9\sim10\); F matches \(10\sim11\), \(9\sim12\). A tour is \(9\to10\to11\to12\to9\). The two separate shapes are separate surfaces; do not combine corners 1--8 with 9--12.}
\hint{1. Choose one edge letter and match arrow tails, then heads. 2. Start a corner card group with one forced pair. 3. When one number occurs in two groups, combine the groups. Keep going until all pairs have been checked.}
\labelled{Important distinction.}{The octagon and square both produce one corner group. That fact alone does not make their local geometry the same. Problem 5 asks how much corner angle belongs to that point.}
\topic{Student p. 5 / Problem 5: how many turns around the point?}
\labelled{Answer and order.}{Arrange octagon sectors in order \(1,6,3,8,5,2,7,4\), closing back to 1, or in reverse cyclic order. Any rotation of this list is equivalent. The sequence of seams is B, C, D, A, B, C, D, A. Square sectors can follow \(9,10,11,12\), closing to 9, with seams E, F, E, F. Rotate the pieces to join corresponding rays; keep both their labels and arrow directions matched.}
\labelled{Angle result.}{The eight \(135^\circ\) sectors give \(1080^\circ\), exactly \textbf{three full turns}. The four \(90^\circ\) sectors give \(360^\circ\), \textbf{one full turn}. Children need not name degrees: an octagon corner is three eighth-turn pieces, so all eight give 24 eighths, or three turns. Each square corner is a quarter-turn.}
\hint{1. First make the ordinary square corner neighborhood. 2. Use that full turn as a reference for the octagon pieces. 3. Count all numbered sectors once, even where they overlap. A circle of visible ink is not the count of angular turns.}
\labelled{Handling and acceptance.}{The p. 5 one-inch reference checks print scale; all pieces are supplied there. Precut if fine-motor work would dominate. Keep the common point fixed, rotate each next sector about it, and maintain the cyclic list as overlapping layers hide pieces. Do not cut away overlap to make the model flat. Accept a correct ordered stack plus a sector-count explanation; a neat flat object is not possible without overlap.}
\labelled{A worthwhile close.}{The square's common corner becomes an ordinary flat point; the octagon's has excess angle. Name the latter a cone point only after children notice the difference. The later L also has three turns, but the current square comparison is a one-square torus, not the three-square L.}
\newpage
\topic{Student p. 6 / Problem 6: six short trips on the L}
\route{Concrete L entry. This is a new map: A/B/C are square names and p/q/r/s are outside edge pairs. Demonstrate one ordinary internal seam and one boundary transfer before independent trials.}
\labelled{Read the actual gluing.}{Crossing right sends A to B across the shared side, B to A across outside pair r, and C to C across pair s. Crossing up sends A to C across the shared side, C to A across pair p, and B to B across pair q. Position along an edge and heading remain unchanged.}
\begin{tabularx}{\linewidth}{@{}p{.65in}XX@{}}
\toprule
Start & Rightward trip & Upward trip \\
\midrule
A & A \(\to\) B \(\to\) A; length 2 & A \(\to\) C \(\to\) A; length 2 \\
B & B \(\to\) A \(\to\) B; length 2 & B \(\to\) B; length 1 \\
C & C \(\to\) C; length 1 & C \(\to\) A \(\to\) C; length 2 \\
\bottomrule
\end{tabularx}\par
\labelled{Complete comparison.}{Lengths use a whole square side as 1. Four trips have length 2: right from A/B and up from A/C. Two have length 1: right from C and up from B. All prescribed dots have local \((1/2,1/4)\); no trip hits a corner. When the route crosses another marked dot, that is not its own return.}
\hint{1. Keep the counter's square name visible. 2. On a boundary jump, compare distance from the arrow tail. 3. Cut a strip one square-side long and compare full paths, including the last piece back to the dot.}
\topic{Student p. 7 / Problem 7: the same dot, a different square}
\labelled{Answer.}{All three starts lie on the same closed orbit, at different places along it. Their successive marked-dot visits are
\[
A\to B\to C\to A,\qquad
B\to C\to A\to B,\qquad
C\to A\to B\to C.
\]
Each complete trip has three whole one-right/one-up blocks, of total length \(3\sqrt2\) square sides. The direction arrow uses a small grid step; it defines a heading, not a turn after every step.}
\labelled{What happens in one block.}{From local \((1/2,1/4)\), right-edge crossing occurs after half a block, top-edge crossing after three-quarters. Apply those two moves in that order. A becomes B then B; B becomes A then C; C becomes C then A. The local position returns at block end, but the square label changes. After three blocks both return.}
\hint{1. Trace one right-and-up block, stopping when local coordinates repeat. 2. Mark which named square contains that dot. 3. Use cards A/B/C to track the block action while the drawn path remains visible.}
\labelled{First return, not merely a return.}{Local coordinates cannot repeat earlier than one whole block, because both local coordinates advance together modulo a square side. The block cycles the three square names, so one or two blocks cannot return to the original state. No simultaneous horizontal/vertical seam event occurs: every prescribed path avoids corners.}
\labelled{Watch for a false shortcut.}{Do not use a square-center dot: a one-right/one-up line from the center hits a vertex. Preserve the supplied starting positions. Reaching the same local position in B after starting in A is a useful observation, not a closed trip.}
\newpage
\topic{Student p. 8 / Problem 8: one heading, two loop lengths}
\labelled{Answer.}{A gives the shortest trip: one two-right/one-up block. B and C each require two blocks and their trips are exactly \textbf{twice as long}. In square-side units the lengths are \(\sqrt5\) and \(2\sqrt5\), but the ratio needs no square roots. B and C lie on the same longer orbit; A lies on a different one.}
\labelled{Exact check of one block.}{At local \((1/4,1/2)\), velocity \((2,1)\), crossings occur right, up, right, at times \(3/8,1/2,7/8\). Apply them in order:
\[
\begin{array}{c|ccc|c}
\text{start}&\text{after right}&\text{after up}&\text{after right}&\text{block-end dot}\\ \hline
A&B&B&A&A\\
B&A&C&C&C\\
C&C&A&B&B.
\end{array}
\]
Thus one block fixes A and swaps B/C. The common local coordinates can return only at integer times, so these are first returns. All event times are distinct; there are no corner hits.}
\hint{1. Compare one full displacement block from each dot. 2. Keep a finger on the named starting square while checking a repeated local dot. 3. Trace the longer loop onto a strip and compare it to two copies of the short one.}
\labelled{A deeper comparison.}{In p. 7, all three marked starts belonged to one orbit. Here they split into a one-block orbit and a two-block orbit. Same heading and local coordinates do not force equal periods when the square label differs.}
\topic{Student p. 9 / Problem 9: four trials, then a general rule}
\route{Readiness-dependent return visit. Allow children to test and revise a conjecture before offering a block argument. Exact calculations below use whole square sides; the printed small-step arrows give the same direction.}
\begin{tabularx}{\linewidth}{@{}p{.65in}p{.84in}X@{}}
\toprule
Right, up & Outcome & Exact evidence from S in A, local \((1/2,1/4)\) \\
\midrule
1, 2 & Closes & One block has up, right, up; its square action fixes A. First return: 1 block, length \(\sqrt5\). \\
2, 3 & Corner stop & At time \(1/4\), the local point is \((1,1)\): A's top-right corner. Stop there, before any ambiguous crossing. Traveled length \(\sqrt{13}/4\). \\
3, 1 & Closes & One block has right, right, up, right. Block-end labels A \(\to\) C \(\to\) B \(\to\) A. First return: 3 blocks, length \(3\sqrt{10}\). \\
3, 2 & Closes & One block has right, up, right, right, up. Block-end labels A \(\to\) B \(\to\) A. First return: 2 blocks, length \(2\sqrt{13}\). \\
\bottomrule
\end{tabularx}\par
\labelled{Universal answer.}{\textbf{Yes}, every positive whole-number right/up direction on this L eventually closes or reaches a corner. If it hits a corner, stop. Otherwise its square-local coordinates repeat after a whole integer displacement. The repeated seam pattern permutes A/B/C, so after at most three such blocks the initial label also returns. Reversibility is essential; the detailed proof is in the adult background.}
\hint{1. Separate a near miss from an exact grid-corner hit. 2. Forget square names briefly and ask when the little position repeats. 3. Restore the square label and follow the effect of that entire block. A record of finitely many experiments alone does not establish the universal answer.}
\newpage
